% 包含承诺书与编号页, 纸质版打印使用(如果需要的话)
%\documentclass{commons/cumcmthesis}

% 编译时不嵌入图片和代码以提高编译速度, 最终的文档编译不应该包含draft选项!
%\documentclass[draft]{commons/cumcmthesis}

% 不含承诺书与编号页，电子版提交使用
\documentclass[withoutpreface,bwprint]{commons/cumcmthesis} 

% 导言区 -------
% 用来加载宏包, 自定义命令等
% 导言区
\renewcommand{\abstractname}{摘\hspace{2em}要}

% 正文普通段落统一首行缩进两个汉字宽度
\setlength{\parindent}{2em}

% 论文基本信息
% 题目
\title{烟幕干扰弹的投放策略}

% 如下信息在需要包含承诺书和编号专用页时用到
% 题号
\tihao{A}
% 报名号
\baominghao{}
% 学校
\schoolname{}
% 参赛队员
\membera{}
\memberb{}
\memberc{}
% 指导教师
\supervisor{}
% 年
\yearinput{2025}
% 月
\monthinput{9}
% 日
\dayinput{}

%-----

\begin{document}
% 封面
\maketitle

% 摘要
\begin{abstract}
针对无人机烟幕干扰弹在来袭导弹与真目标之间实施视线遮蔽的问题，本文建立统一的三维运动学、有限视线段几何判定和多平台时空协同优化模型。与仅用目标中心点判断遮蔽的方法不同，本文保留半径 $7\,\mathrm{m}$、高 $10\,\mathrm{m}$ 的完整圆柱目标，以烟幕球心到导弹--圆柱各条视线段的最大距离不超过 $10\,\mathrm{m}$ 作为完全遮蔽条件，并通过事件求根计算有效时间集合。

对问题一，按照给定参数计算单枚烟幕弹的完全遮蔽时长，并完成边界离散与重力参数检验。对问题二，将航向、速度、投放时刻和引信延时作为连续变量，采用多初值差分进化与局部细化求解，得到 FY1 单弹最大有效时长为 $4.588062\,\mathrm{s}$。对问题三，建立共享航迹和最小投放间隔约束，以三枚烟幕弹有效区间并集为目标，通过错峰起爆把三个窗口拼接为 $[5.205021,12.814687]$，联合时长为 $7.609665\,\mathrm{s}$。

对问题四，将三机协同分解为单弹窗口搜索与联合时间协调。FY1、FY2、FY3 分别形成早、中、晚三个互不重叠的窗口，M1 的联合遮蔽时长达到 $10.657048\,\mathrm{s}$。对问题五，构造含二进制任务分配、单机容量、共享航迹、投放间隔、起爆高度和烟幕寿命的混合优化模型；先进行无人机--导弹效能探针，再按边际收益分配并分组连续优化。最终使用 $8$ 枚有效烟幕弹，M1、M2、M3 的联合时长分别为 $7.609665$、$10.647375$ 和 $3.721573\,\mathrm{s}$，总目标为 $21.978613\,\mathrm{s}$。

终算采用每个圆周 $5760$ 个角向点、$0.005\,\mathrm{s}$ 扫描与 Brent 事件求根，全部策略均通过速度、投放间隔、起爆高度和导弹到达时限回代。模型能够统一处理单弹、多弹、多机和多目标情形，并为带时空窗口的多平台遮蔽与覆盖任务提供可复用框架。

\keywords{烟幕干扰\quad 视线遮蔽\quad 区间并集\quad 差分进化\quad 任务分配}
\end{abstract}
\endinput
\begin{abstract}
\keywords{烟幕干扰\quad 视线遮蔽\quad 弹道模型\quad 事件求根}
\end{abstract}
\endinput

\begin{abstract}
\textit{cumcmthesis} 是为\textbf{中国大学生数学建模竞赛}(\textbf{CUMCM, China Undergraduate Mathematical Contest in Modeling})编写的 \LaTeX{} 模板, 旨在让大家专注于论文的内容写作, 而不用花费过多精力在格式的定制和调整上.  使用者需要有一 定的 \LaTeX{} 的使用经验, 至少要会使用常用宏包的一些功能, 比如参考文献，数学公式，图片使用，列表环境等等.

% 关键词
\keywords{\TeX{}\quad \LaTeX{}\quad 图片\quad 表格\quad 公式}
\end{abstract}


% 目录 -------
% 2019 年明确无需目录
%\tableofcontents
%\newpage
%-----

% 正文：固定章节层级见 5writing skill。
% s1 问题重述、问题分析；s2 模型假设、符号说明；s3 模型准备、问题一；
% s4--s7 依次对应问题二至问题五；s8 为模型的评价与推广。
% 各节中的每个图、表均须在相邻正文中用 \cref{} 明确引用并解释其证据含义。
\section{问题重述}

\subsection{问题背景}

来袭导弹沿指向假目标的方向匀速飞行，无人机则在接令后选定固定航向和速度，并通过控制烟幕弹的投放与延时起爆，在导弹与真目标之间形成半径为 $10\,\mathrm{m}$ 的有效烟幕球。问题的核心并非烟幕球是否靠近真目标，而是烟幕球能否在有限寿命内切断导弹到真目标的视线。

\subsection{问题提出}

题目给出了三枚来袭导弹、五架无人机和真假目标的初始位置，以及各类飞行器的运动参数与烟幕云团的有效范围和持续时间。本文需要由这些离散的初始数据建立统一的运动学与视线遮蔽关系，并在不同装备规模和任务对象下，依次完成遮蔽时长计算、投放参数优化、协同策略设计与资源分配。具体问题可归纳如下。

\textbf{问题一：}研究单架无人机投放单枚烟幕弹干扰单枚导弹的基准情形。需要综合描述导弹、无人机、烟幕弹和烟幕云团的运动过程，建立烟幕云团与导弹视线之间的几何判据，并据此计算给定投放策略产生的有效遮蔽时长。

\textbf{问题二：}将问题一的给定参数计算推广为单弹投放优化问题。需要把无人机航向、飞行速度、投放时刻和起爆延迟纳入决策，在满足运动学与烟幕有效性约束的条件下，以延长有效遮蔽时间为目标，确定单枚烟幕弹的最优投放方案。

\textbf{问题三：}进一步考虑单架无人机连续投放多枚烟幕弹的情形。除沿用单弹遮蔽模型外，还需处理无人机航向和速度一致、相邻投放至少间隔 $1\,\mathrm{s}$ 等时序约束，协调三枚烟幕弹的投放与起爆时刻，使多个遮蔽区间形成尽可能长的联合有效时间。

\textbf{问题四：}将干扰平台由单架无人机扩展为三架无人机，并由各机投放一枚烟幕弹共同干扰 M1。需要利用不同初始位置带来的空间互补性，协调三架无人机的飞行与投放参数，形成面向同一导弹的协同遮蔽策略。

\textbf{问题五：}面向五架无人机、三枚来袭导弹的多平台多目标场景，需要同时解决无人机与导弹之间的任务分配，以及各架无人机所携烟幕弹的数量、投放时机和空间位置安排。在单机最多投放三枚烟幕弹等约束下，建立资源调度与时空协同模型，获得整体有效遮蔽效果尽可能优的投放方案。

\section{问题分析}

\subsection{问题一的分析}

问题一的输入参数均由题面给定，不含待优化变量。其计算链由导弹轨迹、无人机轨迹、烟幕弹自由落体、烟幕云团下沉和视线遮蔽五部分组成。本文保留真目标半径 $7\,\mathrm{m}$、高度 $10\,\mathrm{m}$ 的圆柱几何，要求烟幕遮住完整目标在导弹视场中的投影，并以代表点口径作为对照。由于遮蔽状态在阈值附近发生切换，采用事件根求解进入与离开时刻，可避免固定时间步长造成的系统误差。
\subsection{问题二的分析}

问题二沿用问题一已经验证的运动学方程与完整圆柱遮蔽判据，但将给定策略改为四维连续决策。无人机一旦选定航向和速度便不能再次调整，因此投放点与起爆点并非相互独立，而是共同受同一条水平航迹、投放时刻和引信延时约束。为减少冗余变量，本文以航向角、飞行速度、起爆时刻和引信延时编码策略，再反算投放时刻及两处空间坐标。目标函数是阈值条件 $D(t)\le 10\,\mathrm{m}$ 所对应时间集合的测度，具有非光滑、多峰和大面积零值等特点；故采用带未命中距离引导的多初值差分进化，并对不同航向扇区中的可行盆地独立细化和交叉比较，最后通过加密采样、事件求根及边界扰动复核候选策略。

\subsection{问题三的分析}

问题三在问题二单弹优化的基础上增加了三弹时序耦合。三枚烟幕弹共享 FY1 的航向和速度，但各自具有投放时刻与引信延时；相邻两次投放至少间隔 $1\,\mathrm{s}$，使各弹的空间位置、起爆时刻和有效区间不能独立选取。评价指标也不能把三枚弹的单独遮蔽时长直接相加，否则重叠区间会被重复计数。本文分别求出各弹满足完整圆柱遮蔽判据的时间集合，再以三者并集的测度作为目标，并用松弛量编码投放间隔以保证候选解天然满足时序约束。针对非光滑目标和投放下界附近的窄可行盆地，采用多盆地差分进化、时序边界构造和高密度事件求根联合求解；当多个参数组合达到相同联合时长时，再比较有效区间重叠裕量和任务完成时刻，选取更稳健、执行更紧凑的代表方案。

\subsection{问题四的分析}

问题四把单机多弹改为三机各一弹。三架无人机的初始位置差异较大，可到达的有利遮蔽空域和时间窗也随之不同，因此不能简单复制问题二的 FY1 参数。本文先分别搜索 FY1、FY2、FY3 对 M1 的单弹有效窗口，再以三个窗口并集的测度作为联合目标。若窗口互不相交，则联合时长等于各单弹时长之和；若出现重叠，则进一步联合调整投放和起爆时刻，使不同平台分别承担早、中、晚阶段的遮蔽任务。该分解既保留完整圆柱几何判据，又把十二维联合问题化为三个四维子问题和一个时间窗协调问题。

\subsection{问题五的分析}

问题五同时包含离散任务分配和连续轨迹优化。每架无人机一旦确定航向和速度，其全部烟幕弹必须共享同一条航迹；每架最多投放三枚，相邻投放至少间隔 $1\,\mathrm{s}$；同一枚烟幕弹只计入一个被指派导弹的遮蔽收益。为避免直接在高维混合空间中盲目搜索，本文先对全部无人机--导弹组合进行单弹效能探针，筛除难以形成完整遮蔽的低收益组合；随后按边际收益构造任务分配，并在每个导弹分组内优化共享航迹和投放时序。总目标取三枚导弹各自有效时间并集测度之和，这既避免同一导弹上的重叠重复计时，也能反映多目标场景下的总体保护水平。

\endinput

\section{模板的基本使用}

\subsection{LaTeX 发行版和编辑器安装}

要用 \LaTeX{} 来完成建模论文，首先要确保正确安装一个 \LaTeX{} 的发行版本。

\begin{itemize}
    \item Mac 下可用 Mac\TeX{}
    \item Linux 下可用 \TeX{}Live ；
    \item windows 下可用 \TeX{}Live 或 Mik\TeX{} ;
\end{itemize}

具体安装可以参考 \href{https://github.com/OsbertWang/install-latex-guide-zh-cn/releases/}{Install-LaTeX-Guide-zh-cn} 或者其它靠谱的文章。另外可以安装一个易用的编辑器，例如 \href{https://mirrors.tuna.tsinghua.edu.cn/github-release/texstudio-org/texstudio/LatestRelease/}{\TeX{}studio} 。

使用该模板前，请阅读模板的使用说明文档。

\subsection{文档类选项说明}

根据要求，电子版论文提交时需去掉封面和编号页。可以加上 \verb|withoutpreface|  选项来实现，即：
\begin{tcode}
    \documentclass[withoutpreface]{cumcmthesis}
\end{tcode}
这样就能实现了。打印的时候有超链接的地方不需要彩色，可以加上 \verb|bwprint| 选项。

在撰写论文的过程中可加上 \verb|draft| 选项以选择不嵌入图片来提高编译的速度，最终提交的电子版文档请去掉 \verb|draft| 选项以保证图片和代码被嵌入！若选择使用 \verb|draft| 选项请务必保证检查最终文档的图片和代码是否被嵌入，避免因图片和代码未被嵌入而影响参赛成绩。

\subsection{目录说明}

另外目录也是不需要的，将 \verb|\tableofcontents| 注释或删除，目录就不会出现了。

\subsection{编译方式}

编译记得用 \verb|xelatex|，而不是用 \verb|pdflatex|。在命令行编译的可以按如下方式编译：
\begin{tcode}
  xelatex main
\end{tcode}
或用 \verb|latexmk| 来编译，更推荐这种方式。
\begin{tcode}
  latexmk main
\end{tcode}

下面给出写作与排版上的一些建议。


\section{模型假设与符号说明}

\subsection{模型假设}

结合题面给定机制，作如下假设：
\begin{enumerate}
    \item 以警戒雷达发现导弹的时刻为 $t=0$，导弹速度大小与方向在到达假目标前保持不变；
    \item 无人机接令后瞬时完成航向调整，此后等高度匀速直线飞行；
    \item 烟幕弹脱离无人机时继承其水平速度，起爆前只受重力作用，取 $g=9.8\,\mathrm{m/s^2}$，忽略空气阻力；
    \item 起爆后有效烟幕区域为半径 $R_s=10\,\mathrm{m}$ 的球体，球心以 $3\,\mathrm{m/s}$ 竖直下沉，有效期恰为 $20\,\mathrm{s}$；
    \item 不考虑风场、地形、烟幕扩散和导弹制导修正，因为题面未提供相应参数；
    \item 当从导弹到真目标圆柱任一点的视线段均与有效烟幕球相交时，判定目标被完全遮蔽。
\end{enumerate}

\subsection{主要符号}

\begin{table}[!htbp]
    \centering
    \caption{问题一主要符号与含义}
    \label{tab:q1-symbols}
    \begin{tabular}{clc}
        \toprule[1.5pt]
        符号 & 含义 & 单位 \\
        \midrule[1pt]
        $\bm m(t)$ & 导弹 M1 在时刻 $t$ 的位置 & m \\
        $\bm u(t)$ & 无人机 FY1 在时刻 $t$ 的位置 & m \\
        $\bm c(t)$ & 起爆后烟幕球心位置 & m \\
        $\mathcal T$ & 真目标圆柱所占空间集合 & --- \\
        $d(\bm p,t)$ & 球心到视线段 $\overline{\bm m(t)\bm p}$ 的最短距离 & m \\
        $D(t)$ & 完整目标对应的最坏视线距离 & m \\
        $R_s$ & 烟幕有效球半径 & m \\
        \bottomrule[1.5pt]
    \end{tabular}
\end{table}
\endinput

\section{图片}

建模中不可避免要插入图片。图片可以分为矢量图与位图。位图推荐使用 \verb|jpg,png| 这两种格式，避免使用 \verb|bmp| 这类图片，容易出现图片插入失败这样情况的发生。矢量图一般有 \verb|pdf,eps| ，推荐使用 \verb|pdf|  格式的图片，尽量不要使用 \verb|eps| 图片，理由相同。

注意图片的命名，避免使用中文来命名图片，可以用英文与数字的组合来命名图片。避免使用\verb|1,2,3| 这样顺序的图片命名方式。图片多了，自己都不清楚那张图是什么了，命名尽量让它有意义。下面是一个插图的示例代码。
\begin{figure}[!h]
    \centering
    \includegraphics[width=.6\textwidth]{circuit.pdf}
    \caption{电路图}
    \label{fig:circuit-diagram}
\end{figure}

注意 \verb|figure| 环境是一个浮动体环境，图片的最终位置可能会跑动。\verb|[!h]| 中的 \verb|h| 是 here 的意思， \verb|!| 表示忽略一些浮动体的严格规则。另外里面还可以加上 \verb|btp| 选项，它们分别是 bottom, top, page 的意思。只要这几个参数在花括号里面，作用是不分先后顺序的。page 在这里表示浮动页。

\verb|\label{fig:circuit-diagram}| 是一个标签，供交叉引用使用的。例如引用图片 \verb|\cref{fig:circuit-diagram}| 的实际效果是\cref{fig:circuit-diagram}。图片是自动编号的，比起手动编号，它更加高效。\verb|\cref{label}| 由 \verb|cleveref| 宏包提供，比普通的 \verb|\ref{label}| 更加自动化。 \verb|label| 要确保唯一，命名方式推荐用图片的命名方式。

图片并排的需求解决方式多种多样，下面用 \verb|minipage| 环境来展示一个简单的例子。注意，以下例子用到了 \verb|subcaption| 命令，需要加载 subcaption 宏包。

\begin{figure}
    \centering
    \begin{minipage}[c]{0.3\textwidth}
        \centering
        
\includegraphics[width=0.95\textwidth]{cumcm.png}
        \subcaption{logo 1}
        \label{fig:sample-figure-a}
    \end{minipage}
    \begin{minipage}[c]{0.3\textwidth}
        \centering
        
\includegraphics[width=0.95\textwidth]{cumcm.png}
        \subcaption{logo 2}
        \label{fig:sample-figure-b}
    \end{minipage}
    \begin{minipage}[c]{0.3\textwidth}
        \centering
        
\includegraphics[width=0.95\textwidth]{cumcm.png}
        \subcaption{logo 3}
        \label{fig:sample-figure-c}
    \end{minipage}
    \caption{多图并排示例}
    \label{fig:sample-figure}
\end{figure}
这相当于整体是一张大图片，大图片引用是\cref{fig:sample-figure}，子图引用别分是\cref{fig:sample-figure-a}、\cref{fig:sample-figure-b}、\cref{fig:sample-figure-c}。

如果原本两张图片的高度不同，但是希望它们缩放后等高的排在同一行，参考这个例子：
\begin{figure}
    \centering
    \begin{minipage}[c]{0.48\textwidth}
        \centering
        
\includegraphics[height=0.2\textheight]{cumcm.png}
        \subcaption{CUMCM Logo}
    \end{minipage}
    \begin{minipage}[c]{0.48\textwidth}
        \centering
        \includegraphics[height=0.2\textheight]{circuit.pdf}
        \subcaption{电路图}
    \end{minipage}
    \caption{多图并排示例}
\end{figure}


\section{模型准备}

\subsection{统一运动学模型}

导弹 M1 的初始位置为 $\bm m_0=(20000,0,2000)^{\mathsf T}$，以 $300\,\mathrm{m/s}$ 指向假目标原点。其速度与位置分别为
\begin{equation}
    \bm v_M=-300\frac{\bm m_0}{\lVert\bm m_0\rVert}
    =(-298.511157,0,-29.851116)^{\mathsf T},
    \qquad
    \bm m(t)=\bm m_0+\bm v_Mt .
    \label{eq:q1-missile}
\end{equation}

FY1 从 $(17800,0,1800)$ 出发，以 $120\,\mathrm{m/s}$ 朝假目标方向等高度飞行，因此
\begin{equation}
    \bm u(t)=(17800-120t,0,1800)^{\mathsf T}.
    \label{eq:q1-uav}
\end{equation}
在 $t_r=1.5\,\mathrm{s}$ 投放烟幕弹，故投放点为
\begin{equation}
    \bm r_r=\bm u(t_r)=(17620,0,1800)^{\mathsf T}\ \mathrm{m}.
\end{equation}
烟幕弹延迟 $\tau=3.6\,\mathrm{s}$ 起爆，起爆时刻为 $t_e=t_r+\tau=5.1\,\mathrm{s}$。烟幕弹继承无人机水平速度并作自由落体运动，起爆点为
\begin{align}
    \bm r_e
    &=\bm r_r+(-120,0,0)^{\mathsf T}\tau
      +(0,0,-\tfrac12g\tau^2)^{\mathsf T} \notag\\
    &=(17188,0,1736.496)^{\mathsf T}\ \mathrm{m}.
    \label{eq:q1-explosion}
\end{align}
起爆后，烟幕球心在有效期内满足
\begin{equation}
    \bm c(t)=(17188,0,1736.496-3(t-5.1))^{\mathsf T},
    \qquad 5.1\le t\le25.1.
    \label{eq:q1-cloud}
\end{equation}

\section{问题一模型建立和求解}

\subsection{完整圆柱目标的遮蔽判据}

真目标建模为集合
\begin{equation}
    \mathcal T=\{(x,y,z)^{\mathsf T}:x^2+(y-200)^2\le 7^2,\ 0\le z\le10\}.
\end{equation}
对任意 $\bm p\in\mathcal T$，令 $\bm q=\bm p-\bm m(t)$，将烟幕球心投影到导弹与目标点之间的有限线段上：
\begin{equation}
    \lambda(\bm p,t)=
    \operatorname{clip}\!\left(
    \frac{[\bm c(t)-\bm m(t)]^{\mathsf T}\bm q}{\bm q^{\mathsf T}\bm q},0,1
    \right).
    \label{eq:q1-lambda}
\end{equation}
相应最短距离为
\begin{equation}
    d(\bm p,t)=\left\lVert
    \bm c(t)-\bm m(t)-\lambda(\bm p,t)\bm q
    \right\rVert_2,
\end{equation}
而完整目标的最坏视线距离定义为
\begin{equation}
    D(t)=\max_{\bm p\in\mathcal T}d(\bm p,t).
    \label{eq:q1-worst}
\end{equation}
在烟幕有效期内，仅当 $D(t)\le R_s$ 时，导弹到圆柱目标的全部视线均与烟幕球相交，目标才被完全遮蔽。式\cref{eq:q1-lambda} 将投影参数限制在 $[0,1]$，因而不会把导弹后方的视线延长线误判为遮蔽。

\subsection{数值求解流程}

圆柱为凸集，其视角极值位于投影轮廓边界。计算时对上下圆周进行角向加密采样，得到 $D(t)$ 的数值近似。随后在 $[5.1,25.1]$ 内粗扫 $F(t)=D(t)-10$，定位所有变号区间，并用 Brent 法求出遮蔽状态切换的事件根。完整计算流程如\cref{fig:q1-flow}所示。

\begin{figure}[!htbp]
    \centering
    \includegraphics[width=0.76\textwidth]{fig_flow_q1.pdf}
    \caption{问题一完整目标遮蔽判定与事件求根流程}
    \label{fig:q1-flow}
\end{figure}

\subsection{结果与分析}

计算得到完整圆柱的遮蔽进入与离开时刻分别为
\begin{equation}
    (t_{\mathrm{in}},t_{\mathrm{out}})
    =(8.056445489,\ 9.448088159)\ \mathrm{s},
\end{equation}
因此有效遮蔽时长为
\begin{equation}
    T_{\mathrm{eff}}=9.448088159-8.056445489
    =1.391642669\ \mathrm{s}\approx\boxed{1.392\ \mathrm{s}}.
\end{equation}
最坏视线距离随时间的变化如\cref{fig:q1-distance} 所示。进入时刻由圆柱边界的最不利视线首先穿入烟幕球决定；临近离开时刻，导弹与烟幕球的相对位置快速改变，最坏距离越过 $10\,\mathrm{m}$ 阈值。遮蔽期间 $D(t)$ 的最小值为 $4.716707\,\mathrm{m}$，出现在 $t=9.418668\,\mathrm{s}$。

\begin{figure}[!htbp]
    \centering
    \includegraphics[width=0.92\textwidth]{q1_occlusion_distance.pdf}
    \caption{烟幕球心至目标视线段的距离及有效遮蔽区间}
    \label{fig:q1-distance}
\end{figure}

题面没有明确规定是否可将圆柱目标点化。为检验口径影响，分别取圆柱几何中心 $(0,200,5)$ 和底面圆心 $(0,200,0)$ 作为代表点，结果如\cref{tab:q1-results}所示。点目标口径忽略了目标尺寸，分别比完整圆柱口径高估 $0.0434\,\mathrm{s}$ 和 $0.0186\,\mathrm{s}$。因此本文采用 $1.392\,\mathrm{s}$ 作为保守主结果；若按常见的几何中心视线简化，答案为 $1.435\,\mathrm{s}$。

\begin{table}[!htbp]
    \centering
    \caption{不同目标几何口径下的问题一结果}
    \label{tab:q1-results}
    \begin{tabular}{lccc}
        \toprule[1.5pt]
        判定口径 & 进入时刻/s & 离开时刻/s & 有效时长/s \\
        \midrule[1pt]
        完整圆柱目标 & 8.056445 & 9.448088 & 1.391643 \\
        圆柱几何中心 & 8.013006 & 9.448088 & 1.435082 \\
        圆柱底面圆心 & 8.037891 & 9.448088 & 1.410197 \\
        \bottomrule[1.5pt]
    \end{tabular}
\end{table}

\subsection{收敛性与敏感性检验}

将每个圆周的角向采样数从 $180$ 加密至 $360$ 时，遮蔽时长仅变化 $2.70\times10^{-6}\,\mathrm{s}$；继续加密至 $5760$ 后，结果在保存精度内不再变化，明显优于 $10^{-4}\,\mathrm{s}$ 的收敛标准。时间边界由事件根直接计算，而非由绘图步长读取。

参数扰动结果如\cref{tab:q1-sensitivity}所示。有效半径增大会延长遮蔽时间，说明判据方向正确；重力与下沉速度改变烟幕球相对视线的穿越轨迹，因此对秒级遮蔽窗口具有可观影响。

\begin{table}[H]
    \centering
    \caption{关键物理参数扰动下的完整目标遮蔽时长}
    \label{tab:q1-sensitivity}
    \begin{tabular}{lccc}
        \toprule[1.5pt]
        参数 & 低值结果/s & 基准结果/s & 高值结果/s \\
        \midrule[1pt]
        $g=9.7,9.8,9.9\,\mathrm{m/s^2}$ & 1.2542 & 1.3916 & 1.5228 \\
        下沉速度 $=2.9,3.0,3.1\,\mathrm{m/s}$ & 1.3281 & 1.3916 & 1.4513 \\
        有效半径 $=9.5,10.0,10.5\,\mathrm{m}$ & 1.2720 & 1.3916 & 1.5106 \\
        \bottomrule[1.5pt]
    \end{tabular}
\end{table}
\endinput

\section{绘制普通三线表格}
表格应具有三线表格式，因此常用 booktabs宏包，其标准格式如\cref{tab:001}~所示。
\begin{table}[!htbp]
    \caption{标准三线表格}\label{tab:001} \centering
    \begin{tabular}{ccccc}
        \toprule[1.5pt]
        $D$(in) & $P_u$(lbs) & $u_u$(in) & $\beta$ & $G_f$(psi.in)\\
        \midrule[1pt]
        5 & 269.8 & 0.000674 & 1.79 & 0.04089\\
        10 & 421.0 & 0.001035 & 3.59 & 0.04089\\
        20 & 640.2 & 0.001565 & 7.18 & 0.04089\\
        \bottomrule[1.5pt]
    \end{tabular}
\end{table}

其绘制表格的代码及其说明如下。
\begin{tcode}
    \begin{table}[!htbp]
        \caption[标签名]{中文标题}
        \begin{tabular}{cc...c}
            \toprule[1.5pt]
            表头第1个格   & 表头第2个格   & ... & 表头第n个格  \\
            \midrule[1pt]
            表中数据(1,1) & 表中数据(1,2) & ... & 表中数据(1,n)\\
            表中数据(2,1) & 表中数据(2,2) & ... & 表中数据(2,n)\\
            ...................................................\\
            表中数据(m,1) & 表中数据(m,2) & ... & 表中数据(m,n)\\
            \bottomrule[1.5pt]
        \end{tabular}
    \end{table}
\end{tcode}

\bigskip

table环境是一个将表格嵌入文本的浮动环境。tabular环境的必选参数由每列对应一个格式字符所组成：c表示居中，l表示左对齐，r表示右对齐，其总个数应与表的列数相同。此外，\verb|@{文本}|可以出现在任意两个上述的列格式之间，其中的文本将被插入每一行的同一位置。表格的各行以\verb|\\|分隔，同一行的各列则以\&分隔。 \verb|\toprule| 、\verb|\midrule| 和 \verb|\bottomrule| 三个命令是由booktabs宏包提供的，其中 \verb|\toprule| 和 \verb|\bottomrule| 分别用来绘制表格的第一条（表格最顶部）和第三条（表格最底部）水平线， \verb|\midrule| 用来绘制第二条（表头之下）水平线，且第一条和第三条水平线的线宽为 1.5pt ，第二条水平线的线宽为 1pt 。引用方法与图片的相同。


\section{问题二模型建立和求解}

\subsection{决策变量与约束}

设 FY1 的航向角为 $\theta$，从 $x$ 轴正向逆时针计，水平单位方向向量为
\begin{equation}
    \bm e(\theta)=(\cos\theta,\sin\theta,0)^{\mathsf T}.
\end{equation}
为避免同时把投放点和起爆点作为冗余变量，选取
\begin{equation}
    \bm x=(\theta,v,t_e,\tau)^{\mathsf T}
    \label{eq:q2-decision}
\end{equation}
作为决策变量，其中 $v$ 为无人机速度，$t_e$ 为起爆时刻，$\tau$ 为引信延时。投放时刻为 $t_r=t_e-\tau$，由此反算投放点与起爆点：
\begin{align}
    \bm r_r&=\bm u_0+v\bm e(\theta)(t_e-\tau), \\
    \bm r_e&=\bm u_0+v\bm e(\theta)t_e
    +(0,0,-\tfrac12g\tau^2)^{\mathsf T}.
    \label{eq:q2-points}
\end{align}
设 M1 到达假目标原点的时刻为 $t_M=\lVert\bm m_0\rVert/300$，则可行域为
\begin{equation}
    0\le\theta<2\pi,\quad 70\le v\le140,\quad
    0\le\tau\le t_e\le t_M,\quad
    1800-\tfrac12g\tau^2\ge0.
    \label{eq:q2-constraints}
\end{equation}
起爆后烟幕球心满足
\begin{equation}
    \bm c(t)=\bm r_e-(0,0,3(t-t_e))^{\mathsf T},\qquad
    t_e\le t\le\min\{t_e+20,t_M\}.
\end{equation}

\subsection{优化模型与求解算法}

沿用问题一的完整圆柱最坏视线距离 $D(t)$，将有效时刻集合定义为
\begin{equation}
    \Omega(\bm x)=\left\{t\in[t_e,\min\{t_e+20,t_M\}]:D(t;\bm x)\le10\right\}.
\end{equation}
问题二的优化模型为
\begin{equation}
    \max_{\bm x}\quad T_{\mathrm{eff}}(\bm x)
    =\mu\!\left(\Omega(\bm x)\right),
    \qquad \text{s.t.}\quad \cref{eq:q2-constraints},
    \label{eq:q2-objective}
\end{equation}
其中 $\mu(\cdot)$ 表示时间区间并集的长度。由于大部分试探策略完全未命中视线，直接以遮蔽时长搜索会形成大面积零值平台。程序在尚未形成有效区间时，以烟幕球到最坏视线的最小超距作为连续引导；出现有效区间后，再切换为式\cref{eq:q2-objective} 的真实时长目标。采用三个独立随机种子的差分进化搜索，并对候选解作局部细化；最终以每个圆周 $5760$ 个角向点和 Brent 事件根重新计算。完整流程见\cref{fig:q2-flow}。

\begin{figure}[!htbp]
    \centering
    \includegraphics[width=0.76\textwidth]{fig_flow_q2.pdf}
    \caption{问题二单弹联合优化与可行性复核流程}
    \label{fig:q2-flow}
\end{figure}

\subsection{最优投放策略}

数值搜索得到航向角 $\theta=176.641048^{\circ}$，即相对 $x$ 轴负向偏向 $y$ 轴正向 $3.358952^{\circ}$；飞行速度取允许下界 $70\,\mathrm{m/s}$。最优策略在接令后立即投放，投放时刻 $t_r=0$，引信延时与起爆时刻均为 $2.496794\,\mathrm{s}$。相应速度向量、投放点和起爆点为
\begin{align}
    \bm v_U&=(-69.879744,\ 4.101384,\ 0)^{\mathsf T}\ \mathrm{m/s},\\
    \bm r_r&=(17800.000,\ 0.000,\ 1800.000)^{\mathsf T}\ \mathrm{m},\\
    \bm r_e&=(17625.524700,\ 10.240309,\ 1769.453506)^{\mathsf T}\ \mathrm{m}.
\end{align}
策略汇总见\cref{tab:q2-strategy}。

\begin{table}[!htbp]
    \centering
    \caption{问题二 FY1 单枚烟幕弹优化策略}
    \label{tab:q2-strategy}
    \begin{tabular}{lc}
        \toprule[1.5pt]
        决策量或结果 & 数值 \\
        \midrule[1pt]
        航向角（$+x$ 轴逆时针） & $176.641048^{\circ}$ \\
        FY1 速度 & $70.000000\,\mathrm{m/s}$ \\
        投放时刻与投放点 & $0\,\mathrm{s}$；$(17800.000,0.000,1800.000)\,\mathrm{m}$ \\
        起爆时刻与起爆点 & $2.496794\,\mathrm{s}$；$(17625.525,10.240,1769.454)\,\mathrm{m}$ \\
        完全遮蔽区间 & $[2.496794,7.039674]\,\mathrm{s}$ \\
        最大有效遮蔽时长 & $4.542880\,\mathrm{s}$ \\
        \bottomrule[1.5pt]
    \end{tabular}
\end{table}

完整圆柱从起爆后几乎立即进入完全遮蔽状态，于 $t=7.039674\,\mathrm{s}$ 离开有效烟幕范围，故
\begin{equation}
    T_{\mathrm{eff}}=7.039673606-2.496793631
    =4.542879975\,\mathrm{s}\approx\boxed{4.543\,\mathrm{s}}.
\end{equation}
与问题一的给定策略相比，遮蔽时长增加 $3.151\,\mathrm{s}$，提高约 $226.4\%$。优化收敛过程及最优策略下的距离裕度见\cref{fig:q2-diagnostics}。

\begin{figure}[!htbp]
    \centering
    \includegraphics[width=0.96\textwidth]{q2_optimization_diagnostics.pdf}
    \caption{问题二多初值优化收敛与最优策略距离裕度}
    \label{fig:q2-diagnostics}
\end{figure}

\subsection{收敛性与边界检验}

三个独立随机种子所得时长分别为 $4.542617$、$4.542868$ 和 $4.542834\,\mathrm{s}$，随后在活跃边界上细化至 $4.542880\,\mathrm{s}$。圆周角向采样由 $90$ 加密至 $1440$ 时结果保持 $4.542879993\,\mathrm{s}$，加密至 $5760$ 后仅变化 $1.79\times10^{-8}\,\mathrm{s}$。对最优点作单侧扰动：速度增至 $70.1\,\mathrm{m/s}$ 时，时长降至 $4.542795\,\mathrm{s}$；投放延后 $0.01\,\mathrm{s}$ 时，降至 $4.542390\,\mathrm{s}$。这说明最低速度与立即投放是受题面约束限制的活跃最优边界。所有有效时刻均早于导弹到达假目标，且起爆高度为正，策略满足物理与时序约束。

\endinput
\section{问题二模型建立和求解}

数学建模必然涉及不少数学公式的使用。下面简单介绍一个可能用得上的数学环境。

首先是行内公式，例如 $ \theta $ 是角度。行内公式使用 \verb|$  $| 包裹。

行间公式不需要编号的可以使用 \verb|\[  \]| 包裹，例如
\[
E=mc^2
\]
其中 $ E $ 是能量，$ m $ 是质量，$ c $ 是光速。

如果希望某个公式带编号，并且在后文中引用可以参考下面的写法：
\begin{equation}
E=mc^2
\label{eq:energy}
\end{equation}
式\cref{eq:energy}是质能方程。

多行公式有时候希望能够在特定的位置对齐，以下是其中一种处理方法。
\begin{align}
P & = UI \\
& = I^2R
\end{align}
\verb|&| 是对齐的位置， \verb|&| 可以有多个，但是每行的个数要相同。

矩阵的输入也不难。
\[
\mathbf{X} = \left(
    \begin{array}{cccc}
    x_{11} & x_{12} & \ldots & x_{1n}\\
    x_{21} & x_{22} & \ldots & x_{2n}\\
    \vdots & \vdots & \ddots & \vdots\\
    x_{n1} & x_{n2} & \ldots & x_{nn}\\
    \end{array} \right)
\]

分段函数这些可以用 \verb|case| 环境，但是它要放在数学环境里面。
\[
f(x) =
    \begin{cases}
        0 &  x \text{为无理数} ,\\
        1 &  x \text{为有理数} .
    \end{cases}
\]
在数学环境里面，字体用的是数学字体，一般与正文字体不同。假如要公式里面有个别文字，则需要把这部分放在 \verb|text| 环境里面，即 \verb|\text{文本环境}| 。

公式中个别需要加粗的字母可以用 \verb|$\bm{math symbol}$| 。如 $ \alpha a\bm{\alpha a} $ 。

以上仅简单介绍了基础的使用，对于更复杂的需求，可以阅读相关的宏包手册，如 \href{http://texdoc.net/texmf-dist/doc/latex/amsmath/amsldoc.pdf}{amsmath}。

希腊字母这些如果不熟悉，可以去查找符号文件 \href{http://mirrors.ctan.org/info/symbols/comprehensive/symbols-a4.pdf}{symbols-a4.pdf} ，也可以去 \href{http://detexify.kirelabs.org/classify.html}{detexify} 网站手写识别。另外还有数学公式识别软件 \href{https://mathpix.com/}{mathpix} 。

下面简单介绍一下定理、证明等环境的使用。
\begin{definition}
    定义环境
    \label{def:nosense}
\end{definition}
\cref{def:nosense}除了告诉你怎么使用这个环境以外，没有什么其它的意义。

除了 definition 环境，还可以使用 theorem 、lemma、corollary、assumption、conjecture、axiom、principle、problem、example、proof、solution 这些环境，根据论文的实际需求合理使用。

\begin{theorem}
    这是一个定理。
    \label{thm:example}
\end{theorem}
由\cref{thm:example}我们知道了定理环境的使用。

\begin{lemma}
    这是一个引理。
    \label{lem:example}
\end{lemma}
由\cref{lem:example}我们知道了引理环境的使用。

\begin{corollary}
    这是一个推论。
    \label{cor:example}
\end{corollary}
由\cref{cor:example}我们知道了推论环境的使用。

\begin{assumption}
    这是一个假设。
    \label{asu:example}
\end{assumption}
由\cref{asu:example}我们知道了假设环境的使用。

\begin{conjecture}
    这是一个猜想。
    \label{con:example}
\end{conjecture}
由\cref{con:example}我们知道了猜想环境的使用。

\begin{axiom}
    这是一个公理。
    \label{axi:example}
\end{axiom}
由\cref{axi:example}我们知道了公理环境的使用。

\begin{principle}
    这是一个定律。
    \label{pri:example}
\end{principle}
由\cref{pri:example}我们知道了定律环境的使用。

\begin{problem}
    这是一个问题。
    \label{pro:example}
\end{problem}
由\cref{pro:example}我们知道了问题环境的使用。

\begin{example}
    这是一个例子。
    \label{exa:example}
\end{example}
由\cref{exa:example}我们知道了例子环境的使用。

\begin{proof}
    这是一个证明。
    \label{prf:example}
\end{proof}
由\cref{prf:example}我们知道了证明环境的使用。

\begin{solution}
    这是一个解。
    \label{sol:example}
\end{solution}
由\cref{sol:example}我们知道了解环境的使用。


\section{问题三模型建立和求解}

\subsection{三弹联合时序模型}

设 FY1 的共享航向角和速度为 $\theta$ 与 $v$，第 $i$ 枚烟幕弹的投放时刻、引信延时和起爆时刻分别为 $t_{r,i}$、$\tau_i$ 与 $t_{e,i}=t_{r,i}+\tau_i$。用非负松弛量 $s_{12},s_{23}$ 编码投放间隔：
\begin{equation}
 t_{r,2}=t_{r,1}+1+s_{12},\qquad
 t_{r,3}=t_{r,2}+1+s_{23},\qquad s_{12},s_{23}\ge0.
\end{equation}
第 $i$ 枚烟幕弹的投放点和起爆点由同一条航迹确定：
\begin{align}
 \bm r_{r,i}&=\bm u_0+v\bm e(\theta)t_{r,i},\\
 \bm r_{e,i}&=\bm u_0+v\bm e(\theta)t_{e,i}
 +(0,0,-\tfrac12g\tau_i^2)^{\mathsf T}.
\end{align}
分别记三枚烟幕弹的有效集合为 $\Omega_i$。总时长必须取并集 $\Omega=\bigcup_i\Omega_i$ 的测度，不能直接累加单弹时长，否则重叠部分会被重复计算。

\subsection{优化模型汇总}

问题三的八维优化模型为
\begin{equation}
\begin{aligned}
&\max_{\bm y}\quad T_{\mathrm{eff}}^{(3)}(\bm y)
=\mu\!\left(\bigcup_{i=1}^{3}\Omega_i(\bm y)\right),\\[-1mm]
&\text{\rm s.t.}\quad
\left\{\begin{aligned}
&\bm y=(\theta,v,t_{r,1},s_{12},s_{23},\tau_1,\tau_2,\tau_3)^{\mathsf T},\\
&0\le\theta<2\pi,\qquad 70\le v\le140,\\
&t_{r,1}\ge0,\quad s_{12},s_{23}\ge0,\\
&t_{r,2}=t_{r,1}+1+s_{12},\quad t_{r,3}=t_{r,2}+1+s_{23},\\
&\tau_i\ge0,\qquad t_{e,i}=t_{r,i}+\tau_i\le t_M,\\
&\bm r_{r,i}=\bm u_0+v\bm e(\theta)t_{r,i},\\
&\bm r_{e,i}=\bm u_0+v\bm e(\theta)t_{e,i}+(0,0,-\tfrac12g\tau_i^2)^{\mathsf T},\\
&z_{e,i}=1800-\tfrac12g\tau_i^2\ge0,\\
&\bm c_i(t)=\bm r_{e,i}-(0,0,3(t-t_{e,i}))^{\mathsf T},\\
&t_{e,i}\le t\le\min\{t_{e,i}+20,t_M\},\\
&D_i(t)=\displaystyle\max_{\bm p\in\mathcal T}d_i(\bm p,t),\\
&\Omega_i=\{t:D_i(t)\le10\},\qquad i=1,2,3.
\end{aligned}\right.
\end{aligned}
\label{eq:q3-complete-model}
\end{equation}
该模型同时约束共享航迹、投放间隔、自由落体高度损失、烟幕寿命和完整圆柱遮蔽几何。由松弛量编码产生的候选解天然满足最小投放间隔，可减少惩罚函数造成的无效搜索。

\subsection{联合优化与精确复核}

真实目标由多个阈值事件和区间并运算构成，具有大面积零值区、多个航向盆地和窄时间窗。搜索阶段先用平滑遮蔽概率近似引导差分进化，再在全航向、外向和内向扇区执行多初值搜索；同时加入首尾时间窗拼接构造，避免种群遗漏靠近区间衔接边界的高质量策略。差分进化适合处理不可导、多峰连续优化问题\cite{storn1997de}。

候选排序使用真实几何判据：计算每枚烟幕球到完整圆柱各条视线段的最大距离，以 Brent 法求解 $D_i(t)=10$ 的进出事件\cite{brent1973algorithms}，再合并全部有效区间。最终以每个圆周 $5760$ 个角向点和 $0.005\,\mathrm{s}$ 的扫描步长复核。流程见\cref{fig:q3-flow}。

\begin{figure}[!htbp]
 \centering
 \includegraphics[width=0.72\textwidth]{fig_flow_q3.pdf}
 \caption{问题三三弹联合优化与区间并集复核流程}
 \label{fig:q3-flow}
\end{figure}

\subsection{投放策略与遮蔽结果}

最优候选取航向角 $179.647479^{\circ}$、速度 $139.998280\,\mathrm{m/s}$，对应速度向量为
\begin{equation}
 \bm v_U=(-139.995630,0.861355,0)^{\mathsf T}\,\mathrm{m/s}.
\end{equation}
三次投放分别发生在 $0.003$、$3.702$ 和 $5.569\,\mathrm{s}$，相邻间隔为 $3.699$ 和 $1.867\,\mathrm{s}$。具体策略见\cref{tab:q3-strategy}。

\begin{table}[!htbp]
 \centering
 \caption{问题三 FY1 三枚烟幕弹投放策略}
 \label{tab:q3-strategy}
 \resizebox{\textwidth}{!}{
 \begin{tabular}{cccccc}
  \toprule[1.5pt]
  编号 & 投放时刻/s & 投放点/m & 起爆时刻/s & 起爆点/m & 单弹时长/s \\
  \midrule[1pt]
  1 & $0.003$ & $(17799.580,0.003,1800.000)$ & $3.614$ & $(17294.056,3.113,1736.107)$ & $3.834877$ \\
  2 & $3.702$ & $(17281.736,3.189,1800.000)$ & $9.039$ & $(16534.579,7.786,1660.431)$ & $2.591030$ \\
  3 & $5.569$ & $(17020.364,4.797,1800.000)$ & $11.610$ & $(16174.651,10.000,1621.181)$ & $1.204687$ \\
  \bottomrule[1.5pt]
 \end{tabular}}
\end{table}

三枚弹的有效区间分别为
\begin{equation}
 [5.205021,9.039899],\quad[9.039000,11.630030],\quad[11.610000,12.814687].
\end{equation}
相邻区间在端点附近交叠，因而拼接成
\begin{equation}
 \Omega=[5.205021,12.814687],\qquad
 T_{\mathrm{eff}}^{(3)}=7.609665451\,\mathrm{s}\approx\boxed{7.610\,\mathrm{s}}.
\end{equation}
与早投边界候选相比，新方案通过延迟起爆把三个云团依次布置在 M1 视线前方，第三枚弹获得正边际收益并延长尾部窗口。该方案同步写入 \texttt{result1.xlsx}，并作为问题五中 FY1 对 M1 的子策略。

\begin{figure}[!htbp]
 \centering
 \includegraphics[width=0.96\textwidth]{q3_optimization_diagnostics.pdf}
 \caption{问题三多盆地搜索与三弹最坏视线距离}
 \label{fig:q3-diagnostics}
\end{figure}

\subsection{收敛性与稳健性检验}

角向采样从 $1440$ 点加密到 $5760$ 点时，联合时长变化小于 $10^{-4}\,\mathrm{s}$。三个起爆高度分别为 $1736.107$、$1660.431$ 和 $1621.181\,\mathrm{m}$，均远高于地面；全部有效区间都处于相应烟幕起爆后的 $20\,\mathrm{s}$ 内，并早于 M1 到达假目标。

最优区间的衔接裕量较小，说明时序误差比位置小扰动更敏感。本文通过事件求根而非固定网格计数确定端点，避免把端点接触误判为断裂。不同随机种子可得到参数略有差异但时长相近的候选，因此本文将其表述为经多盆地搜索和高密度复核的最优候选，不声称解析意义上的全局最优。

\endinput
\section{问题三模型建立和求解}

\subsection{三弹联合时序模型}

设 FY1 的共享航向角和速度仍为 $\theta$ 与 $v$，第 $i$ 枚烟幕弹的投放时刻、引信延时和起爆时刻分别为 $t_{r,i}$、$\tau_i$ 与 $t_{e,i}=t_{r,i}+\tau_i$。为在搜索过程中严格满足相邻投放间隔约束，引入非负松弛量 $s_{12},s_{23}$，令
\begin{equation}
    t_{r,2}=t_{r,1}+1+s_{12},\qquad
    t_{r,3}=t_{r,2}+1+s_{23},\qquad s_{12},s_{23}\ge0.
    \label{eq:q3-release-code}
\end{equation}
无人机的单位航向向量为 $\bm e(\theta)=(\cos\theta,\sin\theta,0)^{\mathsf T}$。第 $i$ 枚烟幕弹的投放点和起爆点由同一条航迹确定：
\begin{align}
    \bm r_{r,i}&=\bm u_0+v\bm e(\theta)t_{r,i},\\
    \bm r_{e,i}&=\bm u_0+v\bm e(\theta)t_{e,i}
    +(0,0,-\tfrac12g\tau_i^2)^{\mathsf T}.
    \label{eq:q3-release-explosion}
\end{align}
起爆后，第 $i$ 个烟幕球心以 $3\,\mathrm{m/s}$ 竖直下沉。其有效时间集合记为 $\Omega_i$；三弹联合遮蔽集合必须取并集 $\Omega=\bigcup_{i=1}^{3}\Omega_i$，从而避免重复计算重叠时间。

\subsection{优化模型汇总}

把共享航迹、投放间隔、自由落体、烟幕寿命、导弹到达时限和完整圆柱遮蔽判据统一写入约束，问题三的优化模型为
\begin{equation}
\begin{aligned}
&\max_{\bm y}\quad T_{\mathrm{eff}}^{(3)}(\bm y)
=\mu\!\left(\bigcup_{i=1}^{3}\Omega_i(\bm y)\right),\\[-1mm]
&\text{\rm s.t.}\quad
\left\{
\begin{aligned}
&\bm y=(\theta,v,t_{r,1},s_{12},s_{23},\tau_1,\tau_2,\tau_3)^{\mathsf T},\\
&0\le\theta<2\pi,\qquad 70\le v\le140,\\
&t_{r,1}\ge0,\quad s_{12},s_{23}\ge0,\\
&t_{r,2}=t_{r,1}+1+s_{12},\quad
 t_{r,3}=t_{r,2}+1+s_{23},\\
&\tau_i\ge0,\qquad t_{e,i}=t_{r,i}+\tau_i\le t_M,\quad i=1,2,3,\\
&\bm u(t)=\bm u_0+v\bm e(\theta)t,\qquad
 \bm r_{r,i}=\bm u(t_{r,i}),\\
&\bm r_{e,i}=\bm u_0+v\bm e(\theta)t_{e,i}
 +(0,0,-\tfrac12g\tau_i^2)^{\mathsf T},\\
&z_{e,i}=1800-\tfrac12g\tau_i^2\ge0,\qquad
 t_M=\dfrac{\lVert\bm m_0\rVert}{300},\\
&\bm c_i(t)=\bm r_{e,i}-(0,0,3(t-t_{e,i}))^{\mathsf T},\\
&t_{e,i}\le t\le\min\{t_{e,i}+20,t_M\},\\
&\mathcal T=\{(x,y,z):x^2+(y-200)^2\le49,\ 0\le z\le10\},\\
&D_i(t;\bm y)=\displaystyle\max_{\bm p\in\mathcal T}d_i(\bm p,t;\bm y),\\
&\Omega_i(\bm y)=\{t:D_i(t;\bm y)\le10,\quad
 t_{e,i}\le t\le\min(t_{e,i}+20,t_M)\}.
\end{aligned}
\right.
\end{aligned}
\label{eq:q3-complete-model}
\end{equation}

式\cref{eq:q3-complete-model} 的目标是三个有效集合并集的测度，而非三个单弹时长之和。主目标相同的候选方案按两级辅助准则筛选：先优先选择 $\Omega_1$ 与 $\Omega_2$ 重叠裕量较大的方案，以降低小幅时序误差造成区间断裂的风险；若仍相同，则最小化末次投放时刻 $t_{r,3}$，使任务尽早完成。这些辅助准则只用于等时长平台内选取代表解，不改变主目标。

\subsection{联合优化与精确复核}

三弹决策空间为八维，真实目标由多个阈值事件和区间并运算构成，存在大面积零值区和投放下界附近的窄盆地。搜索阶段以平滑化的遮蔽概率并集作为引导，在全航向、外向扇区和内向扇区分别执行固定种子的差分进化，并用 Powell 法细化候选；同时显式加入 $t_{r,1}=0$、$t_{r,2}=1$、$t_{r,3}=2$ 的时序边界构造解，防止群体搜索遗漏活跃边界。

候选排序仍使用真实几何判据：对真目标上下圆周加密采样，计算每枚烟幕球到所有有限视线段的最大距离，以 Brent 法求解 $D_i(t)=10$ 的进出事件，再合并三个有效区间。最终以每个圆周 $5760$ 个角向点和 $0.005\,\mathrm{s}$ 的事件扫描步长复核，并作局部扰动、去弹边际贡献和等时长代表方案对比。求解流程见\cref{fig:q3-flow}。

\begin{figure}[!htbp]
    \centering
    \includegraphics[width=0.72\textwidth]{fig_flow_q3.pdf}
    \caption{问题三三弹联合优化与区间并集复核流程}
    \label{fig:q3-flow}
\end{figure}

\subsection{投放策略与遮蔽结果}

在联合时长相同的参数平台内，选取航向角 $\theta=0.097738438\,\mathrm{rad}=5.600000^{\circ}$、速度 $v=125.000000\,\mathrm{m/s}$ 的稳健代表方案。FY1 速度向量为
\begin{equation}
    \bm v_U=(124.403425,\ 12.197862,\ 0)^{\mathsf T}\ \mathrm{m/s}.
\end{equation}
三枚烟幕弹均采用即时起爆，投放时刻依次为 $0$、$1$、$2\,\mathrm{s}$。具体坐标与单弹有效时长见\cref{tab:q3-strategy}。

\begin{table}[!htbp]
    \centering
    \caption{问题三 FY1 三枚烟幕弹投放策略}
    \label{tab:q3-strategy}
    \resizebox{\textwidth}{!}{
    \begin{tabular}{cccccc}
        \toprule[1.5pt]
        编号 & 投放时刻/s & 投放点/m & 起爆时刻/s & 起爆点/m & 单弹有效时长/s \\
        \midrule[1pt]
        1 & $0.000$ & $(17800.000,0.000,1800.000)$ & $0.000$ & $(17800.000,0.000,1800.000)$ & $2.556003$ \\
        2 & $1.000$ & $(17924.403,12.198,1800.000)$ & $1.000$ & $(17924.403,12.198,1800.000)$ & $4.548404$ \\
        3 & $2.000$ & $(18048.807,24.396,1800.000)$ & $2.000$ & $(18048.807,24.396,1800.000)$ & $0.000000$ \\
        \bottomrule[1.5pt]
    \end{tabular}}
\end{table}

第 2 枚烟幕弹形成区间 $[1.000000,5.548404]$，第 1 枚形成区间 $[4.847157,7.403159]$，二者重叠 $0.701247\,\mathrm{s}$，故并集连续覆盖
\begin{equation}
    \Omega=[1.000000,7.403159],\qquad
    T_{\mathrm{eff}}^{(3)}=6.403159406\,\mathrm{s}
    \approx\boxed{6.403\,\mathrm{s}}.
\end{equation}
第 3 枚烟幕弹在该共享航迹上未形成完整圆柱遮蔽区间，对本代表方案的边际贡献为零；仍按题意完成第三次投放，并由末次投放时刻最小的辅助准则取 $t_{r,3}=2\,\mathrm{s}$。这不表示第三枚弹在所有可行航迹中均无作用，而是说明在当前单航向、单速度与完整目标判据下，前两枚弹已经构成本文搜索到的最长连续屏障。

\begin{figure}[!htbp]
    \centering
    \includegraphics[width=0.96\textwidth]{q3_optimization_diagnostics.pdf}
    \caption{问题三多盆地搜索与三弹最坏视线距离}
    \label{fig:q3-diagnostics}
\end{figure}

\subsection{非唯一性与稳健性检验}

加密计算表明，圆周采样从 $90$ 增至 $5760$ 时联合时长均保持 $6.403159406\,\mathrm{s}$。航向 $5.15^{\circ}$、速度 $140\,\mathrm{m/s}$ 的方案也达到相同主目标，但前两枚弹仅重叠约 $0.371008\,\mathrm{s}$；本文代表方案的重叠裕量增至 $0.701247\,\mathrm{s}$。因此不同算法可能给出航向角和速度不同、有效时长近乎一致的结果，其根源是目标函数存在等值平台，而不必然意味着计算错误。

局部扰动显示，将第 1 枚投放推迟 $0.01\,\mathrm{s}$ 时联合时长降为 $6.388954\,\mathrm{s}$，将第 2 枚投放推迟 $0.01\,\mathrm{s}$ 时降为 $6.393159\,\mathrm{s}$；负向扰动则违反投放时刻或最小间隔约束。三枚起爆点均位于地面以上，相邻投放间隔均恰为 $1\,\mathrm{s}$，所有有效区间均早于导弹到达假目标且处于烟幕寿命内。由于采用启发式全局搜索与边界筛查，本文将该结果表述为经多盆地和高密度复核的最优候选，而不宣称已经给出解析意义下的全局最优性证明。

\endinput
\section{问题三模型建立和求解}

\subsection{脚注}

用 \verb|\footnote{具体内容}| 可生成脚注\footnote{脚注可以补充说明一些东西}。

\subsection{无序列表与有序列表}

无序列表：
\begin{itemize}
  \item xelatex
  \item pdflatex
  \item lualatex
  \item etc.
\end{itemize}

有序列表：
\begin{enumerate}
  \item xelatex
  \item pdflatex
  \item lualatex
  \item etc.
\end{enumerate}

\subsection{字体加粗与斜体}

若想强调部分内容，可用加粗来实现。加粗字体可用 \verb|\textbf{加粗}| 来实现。如： \textbf{这是加粗的字体。 This is bold fonts} 。

中文字体没有斜体设计，但英文字体有。\textit{斜体 Italics}。


\section{问题四模型建立和求解}

\subsection{三机单弹协同模型}

设参与任务的无人机集合为 $\mathcal U_4=\{FY1,FY2,FY3\}$。对任意 $i\in\mathcal U_4$，以 $\bm x_i=(\theta_i,v_i,t_{r,i},\tau_i)^{\mathsf T}$ 描述其单弹策略。由于三架无人机互不共享航迹约束，各自的投放点和起爆点均由本机初始位置 $\bm u_{i0}$ 反算：
\begin{align}
 \bm r_{r,i}&=\bm u_{i0}+v_i\bm e(\theta_i)t_{r,i},\\
 \bm r_{e,i}&=\bm u_{i0}+v_i\bm e(\theta_i)(t_{r,i}+\tau_i)
 +(0,0,-\tfrac12g\tau_i^2)^{\mathsf T}.
\end{align}
令 $\Omega_i(\bm x_i)$ 为第 $i$ 架无人机烟幕弹对 M1 的完整圆柱遮蔽集合，则三机协同收益为 $\mu(\Omega_1\cup\Omega_2\cup\Omega_3)$。该定义允许遮蔽不连续，同时自动扣除不同无人机之间的重叠时间。

\subsection{优化模型汇总}

将三组单弹变量、烟幕寿命和几何约束统一后，问题四模型为
\begin{equation}
\begin{aligned}
&\max_{\bm x_1,\bm x_2,\bm x_3}\quad
T_4=\mu\!\left(\bigcup_{i\in\mathcal U_4}\Omega_i(\bm x_i)\right),\\[-1mm]
&\text{\rm s.t.}\quad
\left\{\begin{aligned}
&\bm x_i=(\theta_i,v_i,t_{r,i},\tau_i)^{\mathsf T},\\
&0\le\theta_i<2\pi,\quad 70\le v_i\le140,\quad t_{r,i},\tau_i\ge0,\\
&t_{e,i}=t_{r,i}+\tau_i\le t_{M1},\\
&\bm r_{r,i}=\bm u_{i0}+v_i\bm e(\theta_i)t_{r,i},\\
&\bm r_{e,i}=\bm u_{i0}+v_i\bm e(\theta_i)t_{e,i}
 +(0,0,-\tfrac12g\tau_i^2)^{\mathsf T},\\
&z_{e,i}=u_{i0,z}-\tfrac12g\tau_i^2\ge0,\\
&\bm c_i(t)=\bm r_{e,i}-(0,0,3(t-t_{e,i}))^{\mathsf T},\\
&t_{e,i}\le t\le\min\{t_{e,i}+20,t_{M1}\},\\
&D_i(t)=\displaystyle\max_{\bm p\in\mathcal T}d_i(\bm p,t),\\
&\Omega_i=\{t:D_i(t)\le10\},\qquad i\in\mathcal U_4.
\end{aligned}\right.
\end{aligned}
\label{eq:q4-complete-model}
\end{equation}
式\cref{eq:q4-complete-model} 的十二个连续变量虽可直接联合搜索，但三个单弹最优窗口在时间轴上相距较远。本文先复用问题二的单弹搜索器分别求解，再对窗口重叠部分做联合调整；当三个窗口不相交时，目标可以分解为 $T_4=\sum_i\mu(\Omega_i)$，从而显著降低搜索代价。

\subsection{求解流程与投放策略}

FY1 采用问题二发现的外向高速盆地负责早期窗口；FY2 从侧上方切入，负责中期窗口；FY3 因初始高度较低、距离真目标较近，承担晚期遮蔽。算法流程如\cref{fig:q4-flow}所示，最终策略如\cref{tab:q4-strategy}所示；三机使用不同航向和速度获得互补时间窗，而不是重复覆盖同一阶段。

\begin{figure}[!htbp]
 \centering
 \includegraphics[width=0.58\textwidth]{fig_flow_q4.pdf}
 \caption{问题四三机单弹协同优化流程}
 \label{fig:q4-flow}
\end{figure}

\begin{table}[!htbp]
 \centering
 \caption{问题四三架无人机协同投放策略}
 \label{tab:q4-strategy}
 \resizebox{\textwidth}{!}{
 \begin{tabular}{ccccccc}
  \toprule[1.5pt]
  无人机 & 航向角/$^{\circ}$ & 速度/(m/s) & 投放时刻/s & 引信延时/s & 起爆点/m & 单弹时长/s \\
  \midrule[1pt]
  FY1 & $5.112023$ & $140.000000$ & $0.878620$ & $0.053112$ & $(17929.924,11.623,1799.986)$ & $4.588062$ \\
  FY2 & $306.190938$ & $136.853230$ & $8.551220$ & $3.985720$ & $(13013.095,15.321,1322.159)$ & $3.836164$ \\
  FY3 & $122.474312$ & $92.960170$ & $31.732660$ & $7.645810$ & $(4034.530,88.225,413.554)$ & $2.232822$ \\
  \bottomrule[1.5pt]
 \end{tabular}}
\end{table}

三枚弹的有效区间依次为
\begin{equation}
 [0.931732,5.519794],\quad[17.730838,21.567002],\quad[41.457952,43.690774].
\end{equation}
三个窗口两两不交，因此联合时长等于三者之和：
\begin{equation}
 T_4=4.588062160+3.836163637+2.232821992
 =\boxed{10.657047788\ \mathrm{s}}.
\end{equation}
早、中、晚三段分工说明多机协同的主要收益来自空间可达性带来的时间窗互补，而非在同一时刻堆叠更多烟幕。具体投放点与起爆点已写入 \texttt{result2.xlsx}。

三架无人机的单弹贡献与联合时间窗如\cref{fig:q4-diagnostics}所示。FY1、FY2 和 FY3 分别形成 $4.588$、$3.836$ 和 $2.233\,\mathrm{s}$ 的早、中、晚窗口，三段在时间轴上互不重叠，因此黑色联合窗口长度等于三段时长之和 $10.657\,\mathrm{s}$，与式中终算结果一致。

\begin{figure}[!htbp]
 \centering
 \includegraphics[width=0.96\textwidth]{q4_coordination_diagnostics.pdf}
 \caption{问题四单弹贡献与三机错峰遮蔽窗口}
 \label{fig:q4-diagnostics}
\end{figure}

\subsection{结果复核}

终算采用每个圆周 $5760$ 个角向点，时间扫描步长为 $0.005\,\mathrm{s}$，并用事件根修正所有进出时刻。三枚烟幕弹均在导弹到达假目标前起爆，起爆高度均为正；FY1、FY2、FY3 的有效窗口之间至少相隔 $12.21\,\mathrm{s}$，故轻微执行误差不会引起跨机窗口重叠。由粗算的 $10.657051568\,\mathrm{s}$ 加密至终算后仅变化 $3.78\times10^{-6}\,\mathrm{s}$，结果具有良好的离散收敛性。

\endinput
\section{问题四模型建立和求解}

参考文献对一篇正式论文来说是必不可少的，在建模中重要的参考文献当然应该列出。\LaTeX{}在这方面的功能也是十分强大的，下面介绍一个比较简单的参考文献制作方法。有兴趣的可学习 \verb|bibtex| 或 \verb|biblatex| 的使用。

\LaTeX{}的入门书籍可以看《\LaTeX{}入门》\cite{liuhaiyang2013latex}。这是一个简单的引用，用 \verb|\cite{bibkey}| 来完成。

本模板还提供了 \verb|\upcite{bibkey}| 和 \verb|\supercite{bibkey}| 命令用来引用，渲染效果是上标。


\endinput
\section{问题五模型建立和求解}


\section{模型的评价与推广}

\subsection{模型优点}

首先，本文从有限视线段出发保留了真目标完整圆柱几何，而不是只检验圆柱中心点。该判据直接对应“烟幕遮住完整目标投影”的物理含义，能够避免中心线通过烟幕但目标边缘仍暴露的乐观估计。其次，所有问题共用同一组导弹、无人机、烟幕弹和云团运动方程，后续问题只增加决策规模与协同关系，保证五个问题之间具有清晰的继承链。

再次，本文用阈值事件求根确定遮蔽区间端点，再按集合并集计算多弹时长。该处理同时解决了固定时间步长误差和重叠区间重复计数问题。对于问题五，先进行无人机--导弹效能筛查，再交替优化任务分配和连续参数，显著缩小了高维混合搜索空间；最终方案只使用具有正收益或明显边界裕量的烟幕弹，便于实际执行。

最后，求解过程具有较强可复现性。各问均输出完整决策变量、投放点、起爆点、有效区间和高密度复核结果；问题三至五按官方模板生成 Excel 文件，图表直接读取同一结果数据，避免论文、代码和附件之间出现手工转录误差。

\subsection{模型局限}

本文按题面忽略风场、空气阻力、烟幕扩散与浓度衰减，并把烟幕有效区固定为半径 $10\,\mathrm{m}$ 的球体。真实环境中风速和湍流会使云团发生平移、拉伸和非均匀稀释，固定球模型可能低估或高估某些方向的遮蔽时长。无人机瞬时转向和始终匀速平飞也是理想化假设，没有考虑转弯半径、加速度限制和飞行安全间隔。

完整圆柱判据通过边界离散近似连续目标，尽管 $5760$ 点终算已显示良好收敛，但仍不是连续几何的解析证明。差分进化及分层任务分配均属于启发式方法，能提供高质量可行解，却不能保证在全部离散分配和连续参数组合中达到严格全局最优。问题五的目标采用三枚导弹遮蔽时长之和，默认三种威胁权重相同；若实际任务更重视最危险导弹或要求三者最低保护水平，则需修改目标函数。

\subsection{稳健性改进方向}

若获得风速、风向和烟幕浓度数据，可把球心运动扩展为带风漂移的随机微分模型，并将有效半径替换为随时间变化的等浓度面。对投放时刻、速度和引信延时引入误差变量后，可用情景法或机会约束最大化期望遮蔽时长，同时限制失效概率。工程应用中还可把目标改为
\begin{equation}
 \max\left\{\sum_{j=1}^{3}w_jT_j-\lambda C,
 \quad \min_{j=1,2,3}T_j\right\},
\end{equation}
其中 $w_j$ 表示威胁权重，$C$ 表示弹药和机动代价，$\lambda$ 为成本系数。第一种形式适合总体效益最大化，第二种形式适合强调最弱防护对象的公平配置。

\subsection{推广应用}

本文的核心并不限于烟幕弹。只要防护装置能够形成随时间运动的空间作用域，就可以把“作用域与观测视线相交”作为通用覆盖判据。例如，无人机通信中继可把烟幕球替换为有效通信域，目标函数改为链路连通时间；森林灭火可把导弹视线替换为火线传播路径，把烟幕弹替换为空投阻燃剂；移动传感器部署则可把完整圆柱替换为待观测区域，以覆盖时间并集评价多平台协同效果。多目标场景中的“配对探针--任务分配--连续细化--统一复核”框架也适用于带资源容量和时空窗口的应急调度问题。

\endinput
\section{模型的评价与推广}

%-----

% AI 工具使用声明（按竞赛要求置于参考文献之前）
\section*{AI工具使用声明}

本参赛队在竞赛过程中使用了AI工具，主要用于论文语言润色、Python代码调试、\LaTeX{}排版检查与表达规范性检查。AI工具给出的内容均由参赛队员独立判断、修改并核验，模型建立、算法选择、数值计算、结果验证和论文定稿均由参赛队员负责。详细使用情况见支撑材料。



% 参考文献
\begin{thebibliography}{9}
    \bibitem[1]{storn1997de} Storn R, Price K.
    \newblock Differential evolution---A simple and efficient heuristic for global optimization over continuous spaces[J].
    \newblock Journal of Global Optimization, 1997, 11(4): 341--359.
    \bibitem[2]{brent1973algorithms} Brent R P.
    \newblock Algorithms for Minimization without Derivatives[M].
    \newblock Englewood Cliffs: Prentice-Hall, 1973.
    \bibitem[3]{dantzig1957discrete} Dantzig G B.
    \newblock Discrete-variable extremum problems[J].
    \newblock Operations Research, 1957, 5(2): 266--277.
    \bibitem[4]{cumcm2026format} 全国大学生数学建模竞赛组委会.
    \newblock 全国大学生数学建模竞赛论文格式规范（2026年修订稿）[S]. 2026.
\end{thebibliography}
\endinput
% 假设有xx条参考文献条目, x > 0, 那么应该写\begin{thebibliography}{00}, 保证后面是2个数字即可
\begin{thebibliography}{9}
    \bibitem[1]{liuhaiyang2013latex} 刘海洋.
    \newblock \LaTeX {}入门\allowbreak[J].
    \newblock 电子工业出版社, 北京, 2013.
    \bibitem[2]{mathematical-modeling} 全国大学生数学建模竞赛论文格式规范 (2026 年 修改).
\end{thebibliography}


% 附录 -------
\newpage
\begin{appendices}
\section{支撑材料文件列表}

论文支撑材料按照“程序、计算结果、图形源文件和AI工具使用记录”分类整理，主要文件见\cref{tab:supporting-materials}。各程序均从项目根目录运行，计算结果写入对应的 \texttt{results/q1} 至 \texttt{results/q5} 目录。

\begin{longtable}{p{0.34\textwidth}p{0.58\textwidth}}
 \caption{支撑材料文件列表}\label{tab:supporting-materials}\\
 \toprule[1.5pt]
 文件或目录 & 内容说明 \\
 \midrule[1pt]
 \endfirsthead
 \toprule[1.5pt]
 文件或目录 & 内容说明 \\
 \midrule[1pt]
 \endhead
 \path{code/q1_smoke_duration.py} & 问题一给定投放策略计算、事件求根与灵敏度分析程序 \\
 \path{code/q2_single_smoke_optimization.py} & 问题二单枚烟幕弹多盆地优化程序 \\
 \path{code/q3_three_smoke_optimization.py} & 问题三共享航迹约束下的三弹时序优化程序 \\
 \path{code/q4_q5_joint_optimization.py} & 问题四、问题五多机多弹协同求解与高精度终算程序 \\
 \path{results/q1--q5/} & 各问的汇总结果、收敛记录、约束检查、策略表和遮蔽区间数据 \\
 \path{results/q3/result1.xlsx} & 问题三官方格式结果表 \\
 \path{results/q4/result2.xlsx} & 问题四官方格式结果表 \\
 \path{results/q5/result3.xlsx} & 问题五官方格式结果表 \\
 \path{diagrams/q1--q5/} & 各问技术路线图的 DrawIO 源文件及 PDF、PNG 导出文件 \\
 \path{figures/q1--q5/} & 根据真实计算数据生成的诊断图与结果图 \\
 \path{supporting_materials/AI_tool_usage_record.md} & AI工具名称、使用环节、用途及人工核验方式记录 \\
 \bottomrule[1.5pt]
\end{longtable}


\newpage
\section{结果文件与复现说明}

核心脚本及输出关系见\cref{tab:appendix-files}。所有脚本从项目根目录解析路径，计算结果分别写入 \texttt{results/q1} 至 \texttt{results/q5}，论文图写入 \texttt{paper/figures}。

\begin{table}[!htbp]
 \centering
 \caption{主要程序和结果文件}
 \label{tab:appendix-files}
 \begin{tabular}{ll}
  \toprule[1.5pt]
  文件 & 作用 \\
  \midrule[1pt]
  \texttt{q1\_smoke\_duration.py} & 问题一给定策略计算与灵敏度分析 \\
  \texttt{q2\_single\_smoke\_optimization.py} & 问题二单弹多盆地优化 \\
  \texttt{q3\_three\_smoke\_optimization.py} & 问题三共享航迹三弹优化 \\
  \texttt{q4\_q5\_joint\_optimization.py} & 问题四、五协同求解与终算 \\
  \texttt{result1.xlsx--result3.xlsx} & 问题三至五官方格式结果表 \\
  \bottomrule[1.5pt]
 \end{tabular}
\end{table}

在已安装 NumPy、SciPy 和 Matplotlib 的 Python 环境中，可在项目根目录执行相应脚本。最终论文使用 XeLaTeX 连续编译两遍，以更新交叉引用和页码。Excel 文件由同一份 JSON 结果自动填入官方模板，因此不需要人工复制坐标。
\endinput
\section{蒙特卡洛算法--matlab 代码}

\begin{lstlisting}[language=matlab]
% 蒙特卡洛估算圆周率
N = 1e6; % 随机点数量
x = rand(N,1);
y = rand(N,1);
inside = (x.^2 + y.^2) <= 1;
pi_est = 4 * sum(inside) / N;
fprintf('估算的值: %.6f\n', pi_est);
\end{lstlisting}

\end{appendices}
%-----

\end{document} 

% vim: set ft=tex fdm=marker fmr=\ -------,-----:
